

\begin{enumerate}
         \item [c)] La existencia y estado de los equipos de protección radiológica.

    Los equipos detectores de radiación y monitores tipo alarma son los únicos instrumentos con los cuales podemos verificar que la fuente está bajo control. La verificación visual o auditiva de que la fuente se encuentra en el contenedor es muy engañosa, y por lo tanto es muy arriesgado que nos confiemos de ella. El mismo riesgo corremos cuando a los equipos no se les da el mantenimiento o calibración que requieren.

    El monitor tipo alarma sonora es una ayuda valiosa para el técnico, cuando se encuentra ocupado en realizar sus actividades. El incremento de la frecuencia con que emite el sonido le dará idea de cómo se encuentran los niveles de radiación en el lugar en el que se encuentra, sin necesidad de estar continuamente consultando el detector de radiación; sin embargo, sólo el dosímetro personal dará los datos más aproximados a la realidad de la exposición que recibió.

    \item [d)] Condiciones de los sitios en donde se realizaría el trabajo.

    Es innegable que cada lugar en donde se haga la toma de radiografías será distinto y por lo mismo al llegar al área de trabajo, los técnicos deberán quitar de la misma, en lo posible, aquellos objetos que pudieran llegar a ocasionar algún problema al equipo o durante el trabajo.

    El objetivo fundamental de la delimitación de las áreas, de la notificación de que se efectúan trabajos de radiografiado y de la notificación de seguridad es el adecuado manejo de la radiación y la protección de todos los trabajadores involucrados.
\end{enumerate}
